\subsection{Постановка задачи}
\label{sec:analysis:specification}

Требуется разработать автоматизированную информационную систему отслеживания цен на товары. Пользователь взаимодействует с ней через \textit{Telegram}-бота: задаёт критерии товара и пороговую цену, а система опрашивает площадку \textit{Kufar.by}, отбирает подходящие объявления и присылает уведомления о новых предложениях и о снижении цены. Все данные системы хранятся в реляционной базе данных под управлением \textit{PostgreSQL}. В системе предусмотрены две роли. Пользователь отслеживает товары и получает уведомления. Администратор, кроме того, располагает повышенным лимитом отслеживаемых товаров, получает отчёт о состоянии системы и служебные оповещения о сбоях. Администраторы определяются по списку идентификаторов \textit{Telegram}, заданному в конфигурации системы.

\subsubsection{Функциональные требования}

Система хранит сведения о пользователях, их отслеживаемых товарах, объявлениях, полученных с площадки, найденных совпадениях и отправленных уведомлениях, а также справочник характеристик категорий площадки. Пользователь регистрируется автоматически при первом обращении к боту, язык интерфейса (русский или английский) определяется по настройкам \textit{Telegram} и может быть изменён.

Добавление отслеживаемого товара выполняется пошаговым диалогом. Пользователь выбирает категорию, вводит название, выбирает значения характеристик, задаёт пороговую цену и при желании слова-исключения, после чего подтверждает введённые данные. Для категорий «Велосипеды», «Ноутбуки» и «Квартиры (долгосрочная аренда)» отбор ведётся по характеристикам самой площадки: бренду, размеру рамы, объёму памяти, числу комнат, району и другим. В этих категориях у фильтра можно выбрать несколько значений, числовые характеристики (площадь, этаж, год постройки) вводятся диапазоном, а название товара не обязательно. Кроме ручного ввода, пользователь может передать боту ссылку на поиск \textit{Kufar.by}: система переносит из неё категорию, фильтры и цену и сообщает, какие параметры ссылки не поддерживаются. Сразу после добавления товара пользователь получает до пяти самых дешёвых подходящих предложений, уже опубликованных на площадке.

Пользователь просматривает список своих товаров, открывает карточку товара, изменяет его название, цену, слова и фильтры, приостанавливает и возобновляет отслеживание, удаляет товар. Список товаров выгружается в файл формата \textit{CSV} и загружается обратно, причём каждая строка файла проверяется так же, как при ручном вводе.

Система опрашивает площадку каждые 3 минуты и отправляет уведомление в двух случаях: объявление опубликовано после добавления товара и стоит дешевле порога, либо уже известное объявление подешевело. Уведомление содержит фотографии, цену, описание и ссылку на объявление, а также кнопки оценки «Интересно», «Не то» и «Куплю». Одно и то же объявление повторно приходит пользователю только по более низкой цене.

По запросу формируется отчёт в формате \textit{Microsoft Excel} о предложениях, найденных за 7, 30 или 90 дней, с выгодой относительно порога, статусом доставки и оценкой пользователя. Администратору дополнительно доступен отчёт о состоянии системы. Пользователь может привязать адрес электронной почты с подтверждением кодом и получать на него сводные копии уведомлений и отчёты.

Число отслеживаемых товаров ограничено: 3 для пользователя, 10 для администратора и не более 25 для всей системы. Достигнув предела, пользователь получает объяснение, какое именно ограничение сработало.

\subsubsection{Удобство использования}

Отдельного клиентского приложения не требуется: пользователь работает через любой клиент \textit{Telegram} на смартфоне (\textit{Android}, \textit{iOS}), компьютере или в браузере. Управление построено на командах и кнопках под сообщениями, ввод текста нужен только для названия, цены, диапазонов и слов-исключений. Длинные списки значений (бренды, процессоры) разбиваются на страницы по 8 значений. Диалог на любом шаге можно отменить. На каждое действие бот отвечает сообщением: если данные введены неверно, сервис временно недоступен или диалог был прерван, пользователь видит понятное объяснение на своём языке, а не кнопку, которая перестала реагировать.

\subsubsection{Надёжность}

Система должна работать круглосуточно. Каждый компонент выполняется в отдельном контейнере, который при сбое перезапускается автоматически, а состояние контейнеров проверяется регулярными проверками работоспособности. Зависание бота обнаруживается по отсутствию периодической отметки о его активности.

Найденное совпадение сначала сохраняется в базе данных, и только затем отдельная задача отправляет уведомление. Если \textit{Telegram} не принял сообщение, совпадение остаётся в очереди и отправляется повторно, всего до 5 попыток. Уведомления старше 3 суток считаются неактуальными и не отправляются. Сбой почтового сервера не должен задерживать отправку в \textit{Telegram}, поэтому копии уведомлений на электронную почту ведут отдельную очередь.

При отказе площадки обслуживать запросы (ответы с кодами 403 и 429) цикл опроса прерывается, а следующие циклы пропускаются с нарастающей паузой: 1, 2, 4, 8 и 16 циклов. Администратор получает об этом одно оповещение, повторные сообщения на каждом цикле не отправляются. Об остановке, падении или потере работоспособности любого контейнера администратор получает отдельное оповещение.

База данных копируется ежедневно, хранятся 7 последних резервных копий, пригодность копии проверяется восстановлением во временную базу. Структура базы данных изменяется только через версионированные миграции, которые применяются автоматически при запуске системы. Устаревшие данные удаляются по расписанию: объявления без совпадений через 30 дней, записи журнала уведомлений через 90 дней.

\subsubsection{Производительность}

Интервал опроса площадки составляет 3 минуты, к площадке отправляется не более 30 запросов в минуту с паузой между запросами. Каждый отслеживаемый товар требует одного запроса за цикл, и ограничение в 25 товаров на систему гарантирует, что цикл завершается до начала следующего. Если цикл занимает более 80\,\% интервала, система предупреждает об этом в журнале работы. Очередь уведомлений обрабатывается каждую минуту, поэтому от публикации объявления до получения уведомления проходит не более нескольких минут. Чтобы не перегружать пользователя, по одному товару отправляется не более 5 уведомлений в час, а одному пользователю не более 30 в сутки, причём первыми уходят самые дешёвые предложения. Поиск объявления по паре «площадка~-- идентификатор объявления» выполняется на каждом цикле для каждого полученного объявления и обеспечивается составным индексом.

\subsubsection{Безопасность}

Доступ к базе данных имеет только серверный сервис; бот получает и изменяет данные исключительно через его программный интерфейс. Каждый запрос к интерфейсу сопровождается общим секретным ключом, при несовпадении ключа запрос отклоняется, а если ключ не задан, сервис не обслуживает запросы вовсе. Число запросов с одного адреса ограничено 300 в минуту, адрес блокируется на 15 минут после 10 неудачных попыток ввода ключа за 5 минут. Документация программного интерфейса закрыта тем же ключом.

Пользователь видит и изменяет только собственные товары. Код подтверждения адреса электронной почты хранится в виде \textit{HMAC}-хеша, действует 15 минут, допускает 5 попыток ввода и отправляется повторно не чаще раза в минуту. Без подтверждения на адрес ничего не отправляется, поэтому использовать систему для рассылки писем на чужой адрес нельзя. Значения из загружаемых и выгружаемых файлов \textit{CSV}, похожие на формулы, экранируются, а текст объявлений экранируется перед вставкой в сообщение. Секретные параметры (токен бота, пароль базы данных, ключ интерфейса) хранятся в файле окружения вне исходного кода. Контейнеры работают от имени непривилегированного пользователя, а служба мониторинга получает события \textit{Docker} только на чтение через отдельный посредник.

\subsubsection{Аппаратное и программное обеспечение для клиента и сервера}

От клиента требуется только устройство с установленным клиентом \textit{Telegram} или браузером с доступом к веб-версии \textit{Telegram} и подключение к сети Интернет. Для получения копий уведомлений на электронную почту подходит любой почтовый клиент.

Серверная часть рассчитана на один сервер под управлением \textit{Linux} с установленными \textit{Docker} и \textit{Docker Compose}. Рекомендуемая конфигурация сервера: 64-битный процессор с двумя ядрами, не менее 2~ГБ оперативной памяти и не менее 20~ГБ дискового пространства на \textit{SSD} для базы данных, резервных копий и журналов работы. Программная часть реализуется на языке \textit{Python}~3.11 и включает СУБД \textit{PostgreSQL}~16, веб-фреймворк \textit{FastAPI} для серверного интерфейса, библиотеку \textit{aiogram}~3 для работы с \textit{Telegram Bot API}, асинхронный драйвер \textit{asyncpg}, через который выполняются SQL-запросы к базе данных, \textit{Alembic} для версионирования схемы с миграциями в виде SQL-скриптов и \textit{APScheduler} для периодических задач. Для работы системы нужен стабильный доступ к \textit{Kufar.by} и к серверам \textit{Telegram}.
